\documentclass[11pt,a4paper]{article}
\usepackage[margin=25mm]{geometry}
\usepackage{amsmath,amssymb,amsthm,booktabs}
\usepackage[T1]{fontenc}
\usepackage{lmodern}
\usepackage{microtype}
\usepackage[colorlinks=true,linkcolor=blue,urlcolor=blue,citecolor=blue]{hyperref}
\newtheorem{proposition}{Proposition}
\newtheorem{corollary}[proposition]{Corollary}
\newcommand{\Qp}{Q_{+}}
\newcommand{\Qm}{Q_{-}}
\setlength{\parindent}{0pt}
\setlength{\parskip}{6pt}
\title{Prime offsets around even squares:\\local restrictions and finite verification}
\author{Mohammad Aldebbeh}
\date{October 2026}
\begin{document}
\maketitle
\begin{abstract}
For even $n$, we ask whether each adjacent square gap contains a prime
whose distance from $n^2$ is also prime. We prove elementary restrictions,
count admissible offset residues for a fixed center, and compare the two
directions through a simultaneous local count. An exact search, independently
checked by trial division, finds bounded witnesses in both directions for every
even $2\le n\le10\,000$. Published minimal Goldbach-partition computations
give a larger finite range on the lower side. The universal assertions remain
conjectural.
\end{abstract}

\section{The problem}
Let $n\ge2$ be even. An \emph{upper witness} is a prime $q<2n+1$ for
which $n^2+q$ is prime. A \emph{lower witness} is a prime $q<2n-1$ for
which $n^2-q$ is prime. The directions may use different offsets.
Define the bounded minima
\begin{align*}
\Qp(n)&=\min\{q:q\text{ prime},\ q<2n+1,\ n^2+q\text{ prime}\},\\
\Qm(n)&=\min\{q:q\text{ prime},\ q<2n-1,\ n^2-q\text{ prime}\},
\end{align*}
with value $\infty$ when the relevant set is empty. The upper and lower
prime-offset conjectures assert, respectively, that these minima are finite
for every even $n\ge2$.

The strict bounds give
\[
 n^2<n^2+q<(n+1)^2,\qquad (n-1)^2<n^2-q<n^2.
\]
At $n=2$, the minima are $\Qp(2)=3$ and $\Qm(2)=2$, giving primes 7
and 2. For every even $n\ge4$, the offset 2 fails in both directions:
$n^2\pm2$ is even and greater than 2.

\begin{proposition}\label{prop:legendre}
If both prime-offset conjectures hold, every interval
$(m^2,(m+1)^2)$, $m\ge1$ an integer, contains a prime.
\end{proposition}
\begin{proof}
For even $m$, use an upper witness at $n=m$. For odd $m$, use a lower
witness at the even center $n=m+1$. The displayed strict inequalities
place the prime in the required interval. The lower witness at $n=2$
also handles $m=1$.
\end{proof}
This is a conditional implication for Legendre's conjecture. Prime existence
in a square gap by itself imposes no primality condition on the displacement.

\newpage
\section{Elementary arithmetic restrictions}
Recall why coprime integers $a,b$ have integers $x,y$ with $ax+by=1$.
Take the least positive integer $d$ of the form $ax+by$. Division with
remainder of $a$ or $b$ by $d$ shows that any nonzero remainder would
be a smaller positive combination. Thus $d$ divides both $a$ and $b$,
so $d=1$. This is B\'ezout's identity in the coprime case.

We use the elementary fact that a prime $\ell$ dividing $n^2$ divides $n$.
Indeed, if $\ell\nmid n$, B\'ezout's identity gives an inverse of $n$ modulo
$\ell$; multiplying $n^2\equiv0$ by its square would give $1\equiv0$.

\begin{proposition}\label{prop:divisor}
If $n\ge4$ is even and a prime $q$ divides $n$, then both $n^2+q$ and
$n^2-q$ are composite.
\end{proposition}
\begin{proof}
Both values are multiples of $q$. Since $q\le n$,
$n^2-q\ge n^2-n>n\ge q$. Thus each value has a proper prime divisor.
\end{proof}

\begin{corollary}\label{prop:finite}
For every real $B\ge2$, infinitely many even $n$ have no successful prime
offset $q\le B$ in either direction.
\end{corollary}
\begin{proof}
Let $P$ be the product of all primes at most $B$. It includes 2, so all
multiples $n=kP$ are even. For the infinitely many such $n\ge4$, every
prime $q\le B$ divides $n$ and Proposition~\ref{prop:divisor} applies.
\end{proof}
The statement also rules out every fixed finite collection of prime offsets:
choose $B$ at least as large as all its elements. It does not assert that a
larger witness exists. If the two conjectures hold, both minima exceed $B$
at these centers; without them, either minimum could be $\infty$.
For example, $n=30$ eliminates 2, 3 and 5, but its actual minima are 7 and 13.

\begin{proposition}\label{prop:three}
Let $n\ge4$ be even with $3\nmid n$, and let $q>3$ be prime.
Within the respective offset bounds, an upper witness requires
$q\equiv1\pmod3$, and a lower witness requires $q\equiv2\pmod3$.
\end{proposition}
\begin{proof}
Here $n^2\equiv1\pmod3$, while $q$ is congruent to 1 or 2.
If $q\equiv2$, then $3\mid n^2+q$. If $q\equiv1$, then
$3\mid n^2-q$. In the lower bounded interval
$n^2-q>(n-1)^2\ge9$, and the upper value is larger still, so a value
divisible by 3 is composite.
\end{proof}
If $3\mid n$, neither nonzero residue of $q$ is excluded by this modulus.
The offset $q=3$ itself fails by Proposition~\ref{prop:divisor}.
When $3\nmid n$, $q=3$ must instead be tested separately.

More generally, a congruence filter may reject a positive integer $v$
when $\ell\mid v$ and $v>\ell$. It must retain $v=\ell$:
divisibility by a small prime is not sufficient to reject that prime itself.
This equality exception is essential for the lower witness at $n=2$.

\newpage
\section{Counting local conditions at a fixed center}
Fix an even $n$ and a finite set $S$ of distinct primes, with
$M=\prod_{\ell\in S}\ell$. The empty product is 1. For $s\in\{1,-1\}$ set
\[
A_s(n;M)=\{a\bmod M:\gcd(a,M)=\gcd(n^2+sa,M)=1\}.
\]
The variable $a$ ranges over all residue classes, not just primes.
This distinction makes the following count exact.

\begin{proposition}\label{prop:local}
For either sign,
\[
 |A_s(n;M)|=\prod_{\ell\in S}(\ell-\nu_\ell(n)),
 \qquad
 \nu_\ell(n)=
 \begin{cases}1,&\ell\mid n,\\2,&\ell\nmid n.\end{cases}
\]
Moreover, $a\mapsto-a$ is a bijection from $A_+(n;M)$ to $A_-(n;M)$.
\end{proposition}
\begin{proof}
At a single prime $\ell$, exclude $a\equiv0$ and $a\equiv-sn^2$.
These residues coincide exactly when $\ell\mid n^2$, equivalently
$\ell\mid n$. Thus there are $\ell-\nu_\ell(n)$ choices.

For completeness, the Chinese remainder theorem gives a bijection between
residues modulo $M$ and tuples of residues modulo the primes in $S$.
Write $M_\ell=M/\ell$ and choose $u_\ell$ with
$u_\ell M_\ell\equiv1\pmod\ell$, possible by B\'ezout's identity.
The tuple $(a_\ell)$ is represented by
$\sum_{\ell\in S}a_\ell u_\ell M_\ell\bmod M$.
Two representatives of the same tuple differ by a multiple of every prime
in $S$, hence of their product. This proves both existence and uniqueness,
and the local choice counts therefore multiply.

Finally, negation preserves coprimality of $a$ with $M$ and changes
$n^2+a$ into $n^2-(-a)$. Negation is its own inverse.
\end{proof}

The exact fraction in a complete period is
$|A_s|/M=\prod_{\ell\in S}(1-\nu_\ell(n)/\ell)$.
Upper and lower counts agree over a complete residue period even though
their allowed classes can differ. This does not imply equal counts of prime
witnesses in the two bounded intervals.

For an odd $\ell$, the ratio to two independent divisibility tests is
\[
c_\ell(n)=\frac{1-\nu_\ell(n)/\ell}{(1-1/\ell)^2}
=
\begin{cases}
\ell/(\ell-1),&\ell\mid n,\\
\ell(\ell-2)/(\ell-1)^2,&\ell\nmid n.
\end{cases}
\]
At $\ell=2$ the ratio is 2 because $n$ is even. Thus
$C_y(n)=2\prod_{3\le\ell\le y,\ \ell\ {\rm prime}}c_\ell(n)$ is a
finite correction factor for $y\ge2$. Its algebra is exact; interpreting it
as a correction to prime frequencies requires a heuristic model.
The center is fixed throughout: we are counting offset residues, not
solutions of quadratic congruences as $n$ varies.

\newpage
\section{Simultaneous local conditions}
The same offset can be tested on both sides. Put
\[
B(n;M)=A_+(n;M)\cap A_-(n;M).
\]

\begin{proposition}\label{prop:joint}
For even $n$,
\[
 |B(n;M)|=\prod_{\ell\in S}(\ell-\mu_\ell(n)),
 \qquad
 \mu_\ell(n)=
 \begin{cases}
 1,&\ell\mid n,\\
 3,&\ell\nmid n.
 \end{cases}
\]
If $3\in S$ and $3\nmid n$, then $B(n;M)$ is empty.
\end{proposition}
\begin{proof}
The forbidden residues are $0,n^2,-n^2$. If $\ell\mid n$, all three
coincide. Otherwise $\ell$ is odd, since $n$ is even. Then $n^2$ and
$-n^2$ are distinct nonzero residues: equality would imply
$2n^2\equiv0\pmod\ell$. There are consequently three exclusions.
The CRT bijection proved above multiplies the local counts.
For $\ell=3\nmid n$, the factor is $3-3=0$.
\end{proof}

\begin{corollary}\label{prop:common}
For even $n\ge4$ with $3\nmid n$, a prime $q$ that is a bounded witness
in both directions must be 3.
\end{corollary}
\begin{proof}
The offset 2 fails by parity. Every prime $q>3$ has a nonzero residue
modulo 3, and Proposition~\ref{prop:three} requires incompatible residues
for the two directions. Only $q=3$ remains.
\end{proof}
The exception can occur: $4^2-3=13$ and $4^2+3=19$ are both prime.
It is not counted by $B(n;M)$ when $3\mid M$, because the offset itself
equals a sieving prime. At $n=6$, the common witness $q=5$ gives 31 and 41.
There is no contradiction with the two offset conjectures, which allow
different witnesses on the two sides.

The following exact counts use $S=\{2,3,5,7\}$ and $M=210$.
Each entry is checked by enumerating all 210 residues.
\begin{center}
\begin{tabular}{rrrr}
\toprule
$n$ & $|A_+|$ & $|A_-|$ & $|B|$\\
\midrule
% BEGIN LOCAL TABLE
2 & 15 & 15 & 0\\
4 & 15 & 15 & 0\\
6 & 30 & 30 & 16\\
10 & 20 & 20 & 0\\
14 & 18 & 18 & 0\\
30 & 40 & 40 & 32\\
42 & 36 & 36 & 24\\
210 & 48 & 48 & 48\\
% END LOCAL TABLE
\bottomrule
\end{tabular}
\end{center}
These are counts of residue classes. Passing the local tests is not a
certificate of primality. Genuine prime witnesses equal to a sieving prime,
or having a shifted value equal to one, are separate exceptions.

\newpage
\section{Exact bounded computation}
Let $N$ be the largest even integer at most the requested limit.
The reference search marks primes by the Eratosthenes sieve through
$T=(N+1)^2$. It lists prime offsets $2\le q\le2N$ in increasing order.
For each even $2\le n\le N$ and each sign $s$, it tests only
$q<2n+s$, recording the first witness or a failure.

The sieve is exact: every composite $v\le T$ has a prime divisor
$p\le\sqrt v$. Otherwise two nontrivial factors would both exceed
$\sqrt v$. For the least prime divisor, $v\ge p^2$, so the sieve marks
$v$ among the multiples starting at $p^2$. No prime is marked. The
search is in range because the upper shifted value is below
$(n+1)^2\le T$, and the lower one is positive.

The optional filter uses only proper divisors from $\{2,3,5,7\}$; it
retains equality with the divisor. A second calculation uses trial division,
without that filter or the sieve, to check that the recorded offset is prime,
that its shifted value is prime, and that every smaller prime offset fails.
For a reported failure it checks every allowed prime offset.
Trial division is exact for the same square-root divisor reason.

All 5,000 even inputs through $10\,000$ have witnesses in both directions.
The statistics below describe the minima, not all witnesses.
\begin{center}
\begin{tabular}{lrr}
\toprule
Statistic & Upper & Lower\\
\midrule
% BEGIN RESULTS TABLE
Even inputs & $5\,000$ & $5\,000$\\
Failures & 0 & 0\\
Largest minimum & 397 & 467\\
Center of largest minimum & $6\,046$ & $8\,968$\\
Mean & 27.8088 & 22.3494\\
Median & 13 & 11\\
95th percentile & 97 & 83\\
% END RESULTS TABLE
\bottomrule
\end{tabular}
\end{center}
The percentile is the $\lceil0.95k\rceil$-th sorted successful minimum,
where $k=5\,000$. If a search has failures, they are listed explicitly and
descriptive statistics use successful inputs only.

\texttt{results/offsets.csv} retains every input and both minima;
an empty field denotes $\infty$. \texttt{records.csv} contains successive
strict records, \texttt{summary.json} includes the data hash, and
\texttt{local\_counts.csv} contains the residue table. The programs use
only the Python standard library. Tests compare independent searches,
exercise equality and boundary cases, and detect invalid or nonminimal
reported witnesses. Tests supplement the proofs.

The sieve's main array has $(N+1)^2+1$ bytes, about 100 MB at $N=10\,000$,
with additional temporary allocations. The program caps the limit at
20,000. This implementation and these results make no claim for a run to
$10^8$; a different algorithm and retained outputs would be needed.

\newpage
\section{A larger lower-side range from published computation}
Oliveira e Silva, Herzog and Pardi~\cite{goldbach} verified even Goldbach
partitions through $4\cdot10^{18}$. Their $p(E)$ denotes the smallest prime
summand in a partition of the even integer $E>4$.
Section 2.1 and Table 4, printed page 2043, report the successive records of
$p(E)$ over this range. The last and largest record is 9781.
Thus their exhaustive computation supplies, for each such $E$ in range,
primes $r,p$ with
\[
 E=r+p,\qquad r\le9781.
\]
This is an input from a published finite computation, not a theorem
about all even integers and not a computation repeated in this repository.

\begin{proposition}[Using published computational input]\label{prop:goldbach}
Together with the small-case calculation above, that input verifies the
lower prime-offset assertion for every even $2\le n\le2\cdot10^9$.
\end{proposition}
\begin{proof}
If $4892\le n\le2\cdot10^9$ is even, then $n^2\le4\cdot10^{18}$.
Apply the published input at $E=n^2$ and set $q=r$. Then $n^2-q=p$
is prime and
\[
 q\le9781<9783\le2n-1.
\]
The remaining even centers $2\le n\le4890$ are included in the repository's
independently checked search through $10\,000$.
\end{proof}

The threshold is strict: $2(4891)-1=9781$, so the published uniform
bound alone would not suffice there. The first even center above this
threshold is 4892. The even-square endpoint is exact:
$(2\cdot10^9)^2=4\cdot10^{18}$.

This deduction needs the bound on the \emph{minimal} smaller summand.
A Goldbach representation with no bound on its summands would not imply
$q<2n-1$. It also does not yield an upper witness: the upper question asks
for the prime difference $p-q=n^2$, whereas the input gives the sum
$p+q=n^2$. Nor does a prime in the upper square gap guarantee that its
displacement from $n^2$ is prime.

Only two summary facts from the external computation are used: its verified
endpoint and its largest minimal summand. No third-party dataset is
redistributed or silently loaded by the programs. The source links and
page reference are recorded in \texttt{results/README.md}.

The established ranges therefore have different provenance:
\begin{itemize}
\item Both sides through $n=10\,000$: this repository's exact search and
independent verification.
\item Lower side through $n=2\cdot10^9$: an elementary deduction using the
published exhaustive computation and the repository's small cases.
\end{itemize}
No universal prime-existence statement follows from either finite range.

\newpage
\section{Heuristics, limitations and further questions}
A rough model for prime offsets $q<H$ uses approximately $H/\log H$
prime candidates and a prime frequency near $n^2$ of $1/(2\log n)$.
If $H$ is small compared with $n^2$, an uncorrected expected count is
\[
 \lambda_0(n,H)\approx\frac{H}{2\log n\log H}.
\]
Parity doubles this baseline because odd $q$ gives an odd shifted value.
The finite factor $C_y(n)$ describes additional congruence corrections,
but the exact residue counts do not prove that prime candidates in a short
interval follow this model.

Solving the modeled count-of-order-one equation suggests a typical
first-offset scale of order $\log n\log\log n$, up to local factors and
constants. This is a heuristic interpretation, not a proved growth law.
Records over many centers are a different statistic; the finite table gives
no asymptotic bound for either minimum.

For $H$ near $2n$, the model predicts a count of order $n/(\log n)^2$.
A large expected count alone does not force existence. For example,
a random variable equal to 0 or $2n$ with equal probabilities has
expectation $n$ and a fixed probability of zero. A formula such as
$\Pr(\text{no witness})\approx e^{-\lambda}$ requires an extra Poisson
assumption. Even a decreasing modeled failure probability would not
prove success for every even integer.

The proved results concern obstructions, exact finite residue counts, and
a conditional implication between conjectures. The existence and minimum
tables are finite computations. The larger lower range relies on external
computational evidence. The two universal offset assertions, the converse
of Proposition~\ref{prop:legendre}, and asymptotic laws for the minima are
not established here.

Three further investigations follow naturally:
\begin{enumerate}
\item Group the finite minima by $n\bmod30$ or divisibility by small primes.
Compare means and quantiles separately from records.
\item Count common bounded witnesses, separating $3\mid n$ from
$3\nmid n$ and retaining the exceptional offset 3.
\item For a fixed center and cutoff, compare actual witness counts with
locally admissible candidates. Quantify the information lost when a finite
residue filter is used as a model for primality.
\end{enumerate}
None requires enlarging the project's prime-existence claims.

\begin{thebibliography}{9}
\bibitem{goldbach}
T. Oliveira e Silva, S. Herzog and S. Pardi,
\emph{Empirical verification of the even Goldbach conjecture and computation
of prime gaps up to $4\cdot10^{18}$},
Mathematics of Computation \textbf{83} (2014), no.~288, 2033--2060.
\href{https://doi.org/10.1090/S0025-5718-2013-02787-1}{doi:10.1090/S0025-5718-2013-02787-1}.
Section 2.1, Table 4, p.~2043.
\end{thebibliography}

\small
Code and full finite outputs:
\url{https://github.com/Mohammad-Aldebbeh/prime-offsets}.
\end{document}
